\section{付録: プログラムソース}

以下に，実際にコンパイル・実行したプログラムファイルを無加工で添付する．

\subsection{{\tt hashtable\_v1.cpp}（非テンプレート版）}
{\scriptsize\verbatiminput{../src/hashtable_v1.cpp}}

\subsection{{\tt hashtable.cpp}（テンプレート版・提出プログラム）}
{\scriptsize\verbatiminput{../src/hashtable.cpp}}

\subsection{{\tt hashtable\_improved.cpp}（改良版）}
{\scriptsize\verbatiminput{../src/hashtable_improved.cpp}}

\subsection{{\tt exp\_collision.cpp}（実験 1: 衝突分布）}
{\scriptsize\verbatiminput{../src/exp_collision.cpp}}

\subsection{{\tt exp\_anagram.cpp}（実験 2: アナグラム衝突）}
{\scriptsize\verbatiminput{../src/exp_anagram.cpp}}

\subsection{{\tt exp\_loadfactor.cpp}（実験 3: 負荷率と比較回数）}
{\scriptsize\verbatiminput{../src/exp_loadfactor.cpp}}

\subsection{{\tt exp\_timing.cpp}（実験 4: 実行時間）}
{\scriptsize\verbatiminput{../src/exp_timing.cpp}}

\subsection{{\tt exp\_memory.cpp}（実験 5: メモリ使用量）}
{\scriptsize\verbatiminput{../src/exp_memory.cpp}}

\subsection{{\tt exp\_tablesize.cpp}（実験 6: 表サイズの影響）}
{\scriptsize\verbatiminput{../src/exp_tablesize.cpp}}

\subsection{{\tt exp\_rehash.cpp}（実験 7: リハッシュの償却コスト）}
{\scriptsize\verbatiminput{../src/exp_rehash.cpp}}

\subsection{{\tt exp\_openaddr.cpp}（実験 8: 開放番地法との比較）}
{\scriptsize\verbatiminput{../src/exp_openaddr.cpp}}

\subsection{{\tt exp\_keylen.cpp}（実験 9: キー長の影響）}
{\scriptsize\verbatiminput{../src/exp_keylen.cpp}}

\subsection{{\tt test\_hashtable.cpp}（単体テスト）}
{\scriptsize\verbatiminput{../src/test_hashtable.cpp}}

\subsection{{\tt hashtable\_tagged.cpp}（発展版: タグ付き比較）}
{\scriptsize\verbatiminput{../src/hashtable_tagged.cpp}}

\subsection{{\tt exp\_tag.cpp}（実験 10: タグ付き比較の効果）}
{\scriptsize\verbatiminput{../src/exp_tag.cpp}}

\subsection{{\tt exp\_hashdos.cpp}（実験 11: ハッシュ飽和攻撃）}
{\scriptsize\verbatiminput{../src/exp_hashdos.cpp}}

\subsection{{\tt exp\_hashdos.cpp}（実験 11: ハッシュ飽和攻撃の実証）}
{\scriptsize\verbatiminput{../src/exp_hashdos.cpp}}

\subsection{{\tt make\_figures.py}（図の生成スクリプト）}
{\scriptsize\verbatiminput{../experiments/make_figures.py}}
